\documentclass[12pt]{article}
\usepackage[utf8]{inputenc}
\usepackage[T1]{fontenc}
\usepackage{hyperref}
\title{Teke-Tyee Language Resources for K-12 Curriculum Development}
\author{Compiled by Claude Sonnet 5 (claude-sonnet-5) for Andr\'as Kornai}
\date{}
\begin{document}
\maketitle

\section{Introduction}
Teke-Tyee (ISO 639-3 tyx, Glottocode teke1276) is a Bantu language of the Republic of the Congo. Search results report it in the Pool, Bouenza and Lekoumou departments. Speaker counts differ between sources: 14,400 in one and 31,000 in Joshua Project. Other names are Kwe, Tee, Tyee and West Teke.

\section{Primary Educational Resources}
I found no primers, textbooks or school programmes for Teke-Tyee.

\section{Linguistic and Reference Materials}
\begin{itemize}
\item Teke Tyee--French dictionary (SIL, online on Webonary): \url{https://www.webonary.work/teke-tyee/overview/introduction/?lang=en}
\item Dictionary copyright and background page, which describes the Pool, Bouenza and Lekoumou area: \url{https://www.webonary.org/teke-tyee/overview/copyright/?lang=en}
\item West Teke language (Wikipedia): \url{https://en.wikipedia.org/wiki/West_Teke_language}
\item Glottolog entry: \url{https://glottolog.org/resource/languoid/id/teke1276.bigmap.html}
\end{itemize}

\section{Religious and Community Resources}
\begin{itemize}
\item Scripture Earth page for tyx: \url{https://www.scriptureearth.org/00eng.php?iso=tyx}
\item Joshua Project people-group profile: \url{https://joshuaproject.net/people_groups/19589}
\end{itemize}

\section{Digital and Audio Resources}
The Webonary dictionary above is the main digital resource. I did not verify any audio.

\section{Curriculum Development Context}
The dictionary and scripture texts provide a lexical and written base. Primers and reading materials would still have to be created.
\end{document}